%=============================================================================
% NCG GAUGE COUPLING CROSS-REFERENCE:
% Chamseddine--Connes--Marcolli (CCM) Framework vs.\ DFD Modifications
%
% How DFD's CP^2 x S^3 microsector, Toeplitz truncation, traceless
% projection, and regular-module boost modify the standard NCG gauge
% coupling predictions --- and how they change c_1 (the U(1) trace factor)
% and hence kappa_{U(1)}.
%
% Generated: 2026-03-22
%=============================================================================
\documentclass[12pt,a4paper]{article}
\usepackage{amsmath,amssymb,amsthm}
\usepackage{booktabs}
\usepackage{geometry}
\geometry{margin=2.5cm}
\usepackage{hyperref}
\usepackage{tcolorbox}

\newtheorem{theorem}{Theorem}
\newtheorem{finding}{Finding}
\newtheorem{gap}{Gap}

\title{The NCG Gauge Coupling Formula and DFD's Modifications:\\[6pt]
\large From Chamseddine--Connes--Marcolli to the Toeplitz-Truncated Microsector}
\author{Cross-Reference Analysis}
\date{22 March 2026}

\begin{document}
\maketitle

\begin{abstract}
We present a detailed comparison between the standard Chamseddine--Connes--Marcolli
(CCM) noncommutative geometry (NCG) approach to Standard Model gauge couplings
and the Density Field Dynamics (DFD) modifications thereof.  In the CCM framework,
the spectral action on an almost-commutative geometry yields gauge kinetic terms
with coupling ratios $g_2 = g_3 = \sqrt{5/3}\,g_1$ at unification, matching the
$SU(5)$ GUT prediction.  DFD replaces the standard NCG finite space with the
internal geometry $\mathbb{C}P^2 \times S^3$, implements Toeplitz truncation at
$d = 64$, introduces a traceless projection factor $(d-1)/d = 63/64$, and forces
a regular-module boost $\varepsilon_{\mathrm{adj}} = d^2/(d^2-1) = 4096/4095$.
These modifications change the U(1) trace factor $c_1$ and hence $\kappa_{U(1)}$,
yielding $\alpha^{-1} = 137.03599985$ at sub-ppm precision.
\end{abstract}

\tableofcontents
\newpage

%=============================================================================
\section{The Standard CCM Framework}
\label{sec:ccm}
%=============================================================================

\subsection{The Spectral Action Principle}

The Chamseddine--Connes spectral action principle (1996) posits that the
fundamental bosonic action for a noncommutative geometry described by a
spectral triple $(\mathcal{A}, \mathcal{H}, D)$ is:
\begin{equation}\label{eq:spectral-action}
S_{\mathrm{bos}} = \mathrm{Tr}\,f(D_A/\Lambda),
\end{equation}
where $D_A = D + A + \varepsilon JAJ^{-1}$ is the Dirac operator fluctuated by
inner automorphisms (gauge fields), $\Lambda$ is a cutoff scale, and $f$ is a
positive test function.

The asymptotic expansion of this trace in powers of $\Lambda$ yields:
\begin{equation}
S_{\mathrm{bos}} \sim \sum_{k \geq 0} f_k\,\Lambda^{d-2k}\,a_{2k}(D_A^2),
\end{equation}
where $f_k$ are moments of the test function $f$, and $a_{2k}$ are the
Seeley--DeWitt heat-kernel coefficients of $D_A^2$.

\paragraph{Key references.}
\begin{itemize}
\item Chamseddine \& Connes, ``The Spectral Action Principle,'' \textit{Commun.\
Math.\ Phys.}\ \textbf{186}, 731 (1997); arXiv:hep-th/9606001.
\item Chamseddine, Connes \& Marcolli, ``Gravity and the Standard Model with
Neutrino Mixing,'' \textit{Adv.\ Theor.\ Math.\ Phys.}\ \textbf{11}, 991 (2007);
arXiv:hep-th/0610241.
\end{itemize}

\subsection{The Almost-Commutative Geometry of the Standard Model}

The NCG Standard Model uses an almost-commutative geometry:
\begin{equation}
(\mathcal{A}, \mathcal{H}, D) = (\mathcal{A}_M, \mathcal{H}_M, D_M)
\otimes (\mathcal{A}_F, \mathcal{H}_F, D_F),
\end{equation}
where the ``manifold part'' $(\mathcal{A}_M, \mathcal{H}_M, D_M)$ encodes
spacetime and the ``finite part'' $(\mathcal{A}_F, \mathcal{H}_F, D_F)$
encodes the internal gauge and Yukawa structure.

The finite algebra is:
\begin{equation}
\mathcal{A}_F = \mathbb{C} \oplus \mathbb{H} \oplus M_3(\mathbb{C}),
\end{equation}
and the finite Hilbert space $\mathcal{H}_F$ has dimension 32 per generation
(encoding one generation of quarks and leptons with their antiparticles,
chiralities, and color degrees of freedom), giving a total of $32 \times 3 = 96$
complex dimensions for three generations.

\subsection{Gauge Kinetic Terms from the $a_4$ Coefficient}
\label{sec:ccm-gauge}

The gauge kinetic terms arise from the $a_4$ Seeley--DeWitt coefficient in the
heat-kernel expansion.  For a Dirac-type operator $D_A$ on the
almost-commutative space, the $a_4$ term contains:
\begin{equation}\label{eq:gauge-kinetic-ccm}
S_{\mathrm{gauge}} = \frac{f_0}{2\pi^2}
\int d^4x\,\sqrt{g}\left(
c_1\,B_{\mu\nu}B^{\mu\nu}
+ c_2\,W^a_{\mu\nu}W^{a\,\mu\nu}
+ c_3\,G^A_{\mu\nu}G^{A\,\mu\nu}
\right),
\end{equation}
where $B_{\mu\nu}$, $W^a_{\mu\nu}$, $G^A_{\mu\nu}$ are the U(1)$_Y$, SU(2)$_L$,
SU(3)$_c$ field strengths, and $c_1, c_2, c_3$ are trace coefficients computed
over the internal Hilbert space $\mathcal{H}_F$.

\subsection{The Trace Coefficients $c_1, c_2, c_3$}
\label{sec:ccm-traces}

The trace coefficients are computed as:
\begin{equation}
c_r = \mathrm{Tr}_{\mathcal{H}_F}\!\left(T^{(r)}_a T^{(r)}_a\right),
\end{equation}
where $T^{(r)}_a$ are the generators of gauge group $r$ acting on
$\mathcal{H}_F$.

\paragraph{For U(1)$_Y$:}
The hypercharge generator $Y$ acts on each fermion species $f$ with eigenvalue
$Y_f$.  The trace over one generation (with color multiplicity $d_3$ and
weak-isospin multiplicity $d_2$) is:
\begin{equation}
c_1^{(\text{1 gen})} = \sum_f d_3(f)\,d_2(f)\,Y_f^2,
\end{equation}
where the sum runs over the chiral fermion species in one generation.

\paragraph{Explicit computation (per generation, standard $Y$ normalization):}
\begin{center}
\renewcommand{\arraystretch}{1.2}
\begin{tabular}{lcccl}
\toprule
\textbf{Species} & $d_3$ & $d_2$ & $Y$ & $d_3 \cdot d_2 \cdot Y^2$ \\
\midrule
$Q_L = (u_L, d_L)$ & 3 & 2 & $1/6$ & $3 \times 2 \times 1/36 = 1/6$ \\
$u_R$ & 3 & 1 & $2/3$ & $3 \times 1 \times 4/9 = 4/3$ \\
$d_R$ & 3 & 1 & $-1/3$ & $3 \times 1 \times 1/9 = 1/3$ \\
$L_L = (\nu_L, e_L)$ & 1 & 2 & $-1/2$ & $1 \times 2 \times 1/4 = 1/2$ \\
$e_R$ & 1 & 1 & $-1$ & $1 \times 1 \times 1 = 1$ \\
\midrule
\multicolumn{4}{l}{\textbf{Total per generation}} & $\mathbf{10/3}$ \\
\bottomrule
\end{tabular}
\end{center}

With \textbf{GUT normalization} ($Y_{\mathrm{GUT}} = \sqrt{3/5}\,Y$):
\begin{equation}
c_1^{(\text{1 gen, GUT})} = \frac{3}{5} \times \frac{10}{3} = 2.
\end{equation}

\paragraph{For SU(2)$_L$:}
Each SU(2) doublet contributes $\mathrm{Tr}(T^a T^a) = 1/2$ per doublet.
Per generation there are 4 doublets (left-handed quark doublet $\times$ 3 colors
+ lepton doublet):
\begin{equation}
c_2^{(\text{1 gen})} = 4 \times \frac{1}{2} = 2.
\end{equation}

\paragraph{For SU(3)$_c$:}
Each color triplet contributes $\mathrm{Tr}(t^A t^A) = 1/2$ per triplet.
Per generation there are 4 triplets ($Q_L$ has 2 weak components, $u_R$, $d_R$):
\begin{equation}
c_3^{(\text{1 gen})} = 4 \times \frac{1}{2} = 2.
\end{equation}

\paragraph{Three-generation totals:}
\begin{equation}\label{eq:ccm-traces}
\boxed{
c_1 = 3 \times \frac{10}{3} = 10, \qquad
c_2 = 3 \times 2 = 6, \qquad
c_3 = 3 \times 2 = 6.
}
\end{equation}

(With GUT normalization: $c_1^{\mathrm{GUT}} = 3 \times 2 = 6 = c_2 = c_3$.)

\subsection{Gauge Coupling Unification in CCM}
\label{sec:ccm-unification}

Identifying the gauge kinetic terms in Eq.~\eqref{eq:gauge-kinetic-ccm} with
the canonical form $-(1/4g_r^2)\,F^{(r)}_{\mu\nu}F^{(r)\,\mu\nu}$, we obtain:
\begin{equation}\label{eq:ccm-coupling}
\frac{1}{g_r^2} = \frac{2f_0}{\pi^2}\,c_r.
\end{equation}

The unification condition becomes:
\begin{equation}\label{eq:unification}
g_1^2 : g_2^2 : g_3^2 = \frac{1}{c_1} : \frac{1}{c_2} : \frac{1}{c_3}.
\end{equation}

With the three-generation traces from Eq.~\eqref{eq:ccm-traces}:
\begin{equation}
g_1^2 : g_2^2 : g_3^2 = \frac{1}{10} : \frac{1}{6} : \frac{1}{6}.
\end{equation}

This gives the $SU(5)$-type relation:
\begin{equation}\label{eq:su5-relation}
\boxed{g_2 = g_3, \qquad g_1^2 = \frac{3}{5}\,g_2^2,
\qquad\text{i.e.,}\quad g_2 = g_3 = \sqrt{\frac{5}{3}}\,g_1.}
\end{equation}

This is equivalent to $\alpha_3(\Lambda) = \alpha_2(\Lambda)
= (5/3)\,\alpha_1(\Lambda)$ at the unification scale.

\paragraph{The Weinberg angle prediction.}
At unification:
\begin{equation}
\sin^2\theta_W = \frac{g_1^2}{g_1^2 + g_2^2}
= \frac{3/5}{3/5 + 1} = \frac{3}{8}.
\end{equation}
This is the standard GUT prediction $\sin^2\theta_W = 3/8$, which must be
run down to the electroweak scale via renormalization group equations.

\paragraph{The unification problem.}
The CCM framework predicts exact unification at a high scale $\Lambda$.
However, the Standard Model particle content alone does not achieve precise
gauge coupling unification under renormalization group running---the three
couplings do not meet at a single point.  This is a well-known issue shared
with minimal $SU(5)$ GUT models.

%=============================================================================
\section{DFD's Modifications of the NCG Framework}
\label{sec:dfd-mods}
%=============================================================================

DFD departs from the standard CCM framework in four key ways.
Each modification alters how the U(1) trace factor $c_1$ (and hence
$\kappa_{U(1)}$) is computed from the internal geometry.

\subsection{Modification 1: $\mathbb{C}P^2 \times S^3$ Internal Geometry}
\label{sec:mod-geometry}

\paragraph{CCM uses:} An abstract finite spectral triple
$(\mathcal{A}_F, \mathcal{H}_F, D_F)$ with
$\mathcal{A}_F = \mathbb{C} \oplus \mathbb{H} \oplus M_3(\mathbb{C})$.
The internal space is a discrete ``zero-dimensional'' noncommutative space.

\paragraph{DFD uses:} A continuous 7-dimensional internal geometry
$X = \mathbb{C}P^2 \times S^3$.  The gauge structure emerges from the
isometry group of this space:
\begin{itemize}
\item $\mathbb{C}P^2$: isometry group $SU(3)$, providing the $(3,2,1)$
  decomposition that yields $SU(3)_c \times SU(2)_L \times U(1)_Y$.
\item $S^3 \cong SU(2)$: carries the Chern--Simons theory whose partition
  function weights determine the bare lattice couplings.
\end{itemize}

\paragraph{Effect on $c_1$:}
In DFD, the U(1) factor is implemented via twisting by a line bundle over
$\mathbb{C}P^2$ with curvature proportional to the K\"ahler form $\omega$.
The gauge kinetic term emerges from the $a_4$ Seeley--DeWitt coefficient of
the connection Laplacian $P = -g^{ij}\nabla_i\nabla_j$ on this bundle.
The K\"ahler form satisfies $\nabla\omega = 0$, which eliminates derivative
ambiguities in higher heat-kernel coefficients.

The product structure gives a heat-kernel factorization:
\begin{equation}
a_4(X) = a_0(\mathbb{C}P^2) \cdot a_4(S^3)
+ a_2(\mathbb{C}P^2) \cdot a_2(S^3)
+ a_4(\mathbb{C}P^2) \cdot a_0(S^3),
\end{equation}
and the gauge kinetic term arises from the
$\Omega_{\mu\nu}\Omega^{\mu\nu}$ piece of $a_4(\mathbb{C}P^2)$.

\subsection{Modification 2: Toeplitz Truncation at $d = 64$}
\label{sec:mod-toeplitz}

\paragraph{CCM uses:} The full infinite-dimensional function algebra on the
internal space (or equivalently, the finite-dimensional algebra
$\mathcal{A}_F$ directly).

\paragraph{DFD uses:} Toeplitz quantization at level $m = k + 3$ on
$\mathbb{C}P^1 \subset \mathbb{C}P^2$, replacing the continuous geometry
with a finite matrix algebra $M_d(\mathbb{C})$ where:
\begin{equation}
d = \dim H^0(\mathbb{C}P^1, \mathcal{O}(m)) = m + 1 = k + 4 = 64
\quad (k_{\max} = 60).
\end{equation}

The Spin$^c$ shift $+3$ arises from the canonical bundle
$K_{\mathbb{C}P^2} = \mathcal{O}(-3)$, giving
$L_{\det} = K^{-1} = \mathcal{O}(3)$.  On restriction to
$\mathbb{C}P^1$, the bundle $\mathcal{O}(k) \otimes L_{\det}$ becomes
$\mathcal{O}(k+3)$.

The value $k_{\max} = 60$ is derived from the holomorphic Euler
characteristic:
\begin{equation}
k_{\max} = \chi(\mathbb{C}P^2, \mathcal{O}(9) \oplus \mathcal{O}^{\oplus 5})
= \binom{11}{2} + 5 = 55 + 5 = 60.
\end{equation}

\paragraph{Effect on $c_1$:}
The Toeplitz truncation replaces the continuous integral in $a_4$ with a
finite-dimensional matrix trace over the $d$-dimensional truncated Hilbert
space.  This modifies the overall normalization of $c_1$ by replacing the
infinite spectral sum with a sum over at most $d = 64$ modes.  The effective
spectral cutoff becomes $\Lambda^3 = k(d-1)/d = 60 \times 63/64 = 59.0625$.

\subsection{Modification 3: Traceless Projection $(d-1)/d = 63/64$}
\label{sec:mod-traceless}

\paragraph{CCM:} The unimodularity condition removes the overall U(1) from the
gauge group, but this is implemented at the algebraic level (as a constraint
on the algebra $\mathcal{A}_F$) rather than as a multiplicative correction to
the trace.

\paragraph{DFD:} The Toeplitz-truncated algebra $M_d(\mathbb{C})$ has gauge
group $\mathrm{U}(d) = \mathrm{U}(1) \times \mathrm{SU}(d)/\mathbb{Z}_d$.
The determinant (trace) channel is removed explicitly, projecting onto the
traceless part.  This produces the multiplicative factor:
\begin{equation}
\boxed{\text{Traceless projection} = \frac{d-1}{d} = \frac{63}{64} = 0.984375.}
\end{equation}

\paragraph{Effect on $c_1$:}
The traceless projection reduces the effective number of modes contributing
to the U(1) gauge kinetic term by the factor $63/64$.  In the CCM framework,
this factor is unity (or absorbed into the algebraic structure).  In DFD, it
enters the spectral cutoff $\Lambda^3 = k(d-1)/d$ and hence modifies
$\kappa_{U(1)}$ multiplicatively.

\subsection{Modification 4: Regular-Module Boost
$\varepsilon_{\mathrm{adj}} = d^2/(d^2-1) = 4096/4095$}
\label{sec:mod-boost}

\paragraph{CCM:} The finite Hilbert space $\mathcal{H}_F$ is a
\emph{bimodule} over $\mathcal{A}_F$, carrying the fermion representation.
Traces are computed in the fermion representation directly.

\paragraph{DFD:} The finite Hilbert space is taken as the algebra itself
(the ``regular module''):
\begin{equation}
\mathcal{H}_F := A = M_d(\mathbb{C}), \qquad \dim(\mathcal{H}_F) = d^2 = 4096.
\end{equation}

The natural UV-normalized (democratic) trace averages over all $d^2 = 4096$
matrix directions:
\begin{equation}
\mathrm{tr}_{\mathrm{dem}}(\cdot) = \frac{1}{d^2}\,\mathrm{Tr}_{\mathcal{H}_F}(\cdot).
\end{equation}

But physical gauge fields live in $\mathfrak{su}(d)$ (dimension $d^2 - 1 = 4095$),
not $\mathfrak{u}(d)$ (dimension $d^2 = 4096$).  Converting from democratic
to physics normalization forces:
\begin{equation}
\boxed{\varepsilon_{\mathrm{adj}}^{(A)} = \frac{d^2}{d^2-1} = \frac{4096}{4095}
= 1.000244\ldots \qquad\text{(BOOST)}}
\end{equation}

This is the ``forced binary fork'': the only alternative (Branch B,
$\varepsilon^{(B)} = 4095/4096$) is falsified at 43~ppm against the
experimental $\alpha^{-1}$.

\paragraph{Effect on $c_1$:}
The BOOST factor multiplies $\kappa_{U(1)}$ by $4096/4095$, shifting the
predicted $\alpha^{-1}$ by $\sim 0.03$ (from $\sim 137.003$ to $137.036$).
In the CCM framework, no such factor appears because the fermion Hilbert
space is not the regular module.

%=============================================================================
\section{How DFD Changes $c_1$ and $\kappa_{U(1)}$}
\label{sec:c1-change}
%=============================================================================

\subsection{The CCM Prediction for $\kappa_{U(1)}$}

In the standard CCM framework, the gauge coupling relation
(Eq.~\ref{eq:ccm-coupling}) gives:
\begin{equation}
\frac{1}{g_1^2} = \frac{2f_0}{\pi^2}\,c_1 = \frac{2f_0}{\pi^2} \times 10
\qquad (\text{3-generation, standard } Y).
\end{equation}

The stiffness $\kappa_{U(1)} = 1/g_1^2$ is then determined entirely by the
moment $f_0$ and the trace $c_1 = 10$.  The CCM framework predicts
\emph{unification} ($g_1^2 = (3/5)g_2^2 = (3/5)g_3^2$) at some scale
$\Lambda$, but does not directly predict the low-energy value of $\alpha$.

\subsection{The DFD Prediction for $\kappa_{U(1)}$}

DFD replaces the CCM formula with a modified computation that incorporates
the four modifications above.  The DFD prediction is structured as:
\begin{equation}\label{eq:kappa-dfd}
\boxed{
\kappa_{U(1)} = \kappa_{\mathrm{raw}} \times w
\times \varepsilon_{\mathrm{adj}} \times \frac{k+3}{k+4},
}
\end{equation}
where:
\begin{itemize}
\item $\kappa_{\mathrm{raw}}$ is the universal gauge-kinetic prefactor from
  the $a_4$ Seeley--DeWitt coefficient on the Toeplitz-truncated
  $\mathbb{C}P^2 \times S^3$,
\item $w = N_{\mathrm{species}}/(g_F \cdot \mathrm{Tr}(Y^2)) = 7/80$ is the
  hypercharge weighting that converts the universal coefficient to the U(1)
  sector,
\item $\varepsilon_{\mathrm{adj}} = d^2/(d^2-1) = 4096/4095$ is the
  regular-module boost,
\item $(k+3)/(k+4) = 63/64$ is the traceless projection factor.
\end{itemize}

\subsection{Comparison: CCM $c_1$ vs.\ DFD $c_1$}

\begin{center}
\renewcommand{\arraystretch}{1.4}
\begin{tabular}{p{4cm}cc}
\toprule
\textbf{Feature} & \textbf{CCM} & \textbf{DFD} \\
\midrule
Internal geometry & Discrete finite triple & $\mathbb{C}P^2 \times S^3$ \\
$c_1$ (standard $Y$, 3 gen) & $10$ & Absorbed into $\kappa_{\mathrm{raw}} \times w$ \\
$c_1$ (GUT-normalized, 3 gen) & $6$ & Not used \\
Traceless projection & $1$ (algebraic) & $63/64$ \\
Regular-module boost & $1$ (no boost) & $4096/4095$ \\
Spectral truncation & None (finite algebra) & Toeplitz at $d = 64$ \\
$\mathrm{Tr}(Y^2)$ & $10/3$ per gen & $10$ (3-gen sum) \\
$g_F$ (fermion multiplicity) & Implicit in $\dim\mathcal{H}_F = 96$ & $8$ (explicit) \\
$N_{\mathrm{species}}$ & Implicit & $7$ (SU(2) components) \\
\bottomrule
\end{tabular}
\end{center}

The key difference is that DFD \textbf{decomposes} $c_1$ into explicit
geometric and SM-content factors ($w$, $\varepsilon_{\mathrm{adj}}$,
$(d-1)/d$), whereas CCM treats $c_1$ as a single trace over the finite
Hilbert space.

%=============================================================================
\section{The Electroweak Mixing and the $\alpha^{-1}$ Formula}
\label{sec:alpha}
%=============================================================================

\subsection{CCM Route to $\alpha$}

In the CCM framework, at the unification scale:
\begin{equation}
\frac{1}{e^2} = \frac{1}{g_1^2} + \frac{1}{g_2^2}
= \frac{2f_0}{\pi^2}(c_1 + c_2)
= \frac{2f_0}{\pi^2}(10 + 6) = \frac{32 f_0}{\pi^2}.
\end{equation}

This does not directly predict $\alpha$ at low energies; it requires running
the couplings down from the unification scale $\Lambda$ to $m_Z$.

\subsection{DFD Route to $\alpha$}

DFD uses the Frame Stiffness Theorem (from the Ricci curvature of the
$(3,2,1)$ partition on $\mathbb{C}P^2$):
\begin{equation}
\kappa_r = n_r \kappa_0, \qquad n_{U(1)} = 1, \quad n_{SU(2)} = 2,
\quad n_{SU(3)} = 3.
\end{equation}

The gauge couplings are extracted via Wilson normalization:
\begin{equation}
g_1^2 = \frac{1}{\kappa_{U(1)}}, \qquad
g_2^2 = \frac{4}{\kappa_{SU(2)}}.
\end{equation}

Electroweak mixing with $\kappa_{SU(2)} = 2\kappa_{U(1)}$:
\begin{equation}
\frac{1}{e^2} = \kappa_{U(1)} + \frac{2\kappa_{U(1)}}{4}
= \frac{3}{2}\,\kappa_{U(1)}.
\end{equation}

Therefore:
\begin{equation}
\boxed{\alpha^{-1} = \frac{4\pi}{e^2} = 6\pi\,\kappa_{U(1)}.}
\end{equation}

The DFD prediction is $\kappa_{U(1)} = 7.26999\ldots$, giving:
\begin{equation}
\alpha^{-1} = 6\pi \times 7.26999\ldots = 137.03599985
\qquad (\text{residual: } {-0.006}\text{ ppm}).
\end{equation}

\paragraph{The Wilson ratio.}
The lattice coupling ratio is
$\beta_{SU(2)}/\beta_{U(1)} = (n_2/n_1) \times N_{\mathrm{gen}} = 2 \times 3 = 6$.

\subsection{Critical Difference: DFD Does Not Require Unification}

The CCM framework requires all couplings to unify at a single scale
$\Lambda$, which is problematic because the SM running couplings do not
meet at one point.  DFD sidesteps this by:
\begin{enumerate}
\item Computing $\kappa_{U(1)}$ directly from the Chern--Simons weighted
  average on the microsector (no unification scale needed).
\item Deriving the stiffness \emph{ratios} from frame geometry
  ($\kappa_r \propto n_r$) rather than from a single $f_0$ coefficient.
\item Predicting $\alpha$ at \emph{low energy} directly, without
  renormalization group running from a GUT scale.
\end{enumerate}

%=============================================================================
\section{The Bare Coupling $\beta_{U(1)}$ from Chern--Simons Theory}
\label{sec:cs}
%=============================================================================

The DFD microsector places an $SU(2)_k$ Chern--Simons theory on the $S^3$
factor.  The Euclidean weight at CS level $k$ is:
\begin{equation}
w(k) = |Z_{SU(2)_k}(S^3)|^2 = \frac{2}{k+2}\,\sin^2\!\frac{\pi}{k+2},
\end{equation}
and the bare U(1) lattice coupling is the weighted average of shifted levels:
\begin{equation}
\beta_{U(1)} = \langle k_{\mathrm{eff}} \rangle
= \frac{\sum_{k=0}^{59}(k+2)\,w(k)}{\sum_{k=0}^{59}w(k)}
= 3.7969\ldots \approx 3.80.
\end{equation}

In the CCM framework, no such Chern--Simons computation exists.  The gauge
couplings are determined by the single moment $f_0$ and the trace $c_r$.
DFD replaces $f_0$ with a \emph{computable} quantity: the CS partition function
weighted average.

%=============================================================================
\section{Summary: CCM vs.\ DFD Gauge Coupling Derivation}
\label{sec:summary}
%=============================================================================

\begin{tcolorbox}[colback=blue!5!white, colframe=blue!50!black,
  title=CCM vs.\ DFD: Gauge Coupling Comparison]

\textbf{Standard CCM (Chamseddine--Connes--Marcolli):}
\begin{itemize}
\item Gauge kinetic term: $\propto (f_0/2\pi^2)\,c_r\,F_r^2$
\item Trace coefficients: $c_1 = 10$, $c_2 = c_3 = 6$ (3 generations)
\item Coupling ratios: $g_2 = g_3 = \sqrt{5/3}\,g_1$ (SU(5)-type)
\item Unification: required at scale $\Lambda$; Weinberg angle $\sin^2\theta_W = 3/8$
\item Free parameter: $f_0$ (moment of test function)
\item Low-energy $\alpha$: requires RG running, \textbf{not directly predicted}
\end{itemize}

\textbf{DFD Modifications:}
\begin{enumerate}
\item \textbf{Internal geometry:} $\mathbb{C}P^2 \times S^3$ replaces
  discrete finite triple
\item \textbf{Toeplitz truncation:} $d = 64$ replaces infinite function space;
  modifies $a_4$ normalization
\item \textbf{Traceless projection:} $(d-1)/d = 63/64$ from $SU(d)$ vs.\ $U(d)$
  distinction
\item \textbf{Regular-module boost:} $\varepsilon_{\mathrm{adj}} = 4096/4095$
  from regular module $\mathcal{H}_F = M_d(\mathbb{C})$
\end{enumerate}

\textbf{DFD Result:}
\begin{itemize}
\item $\kappa_{U(1)} = \kappa_{\mathrm{raw}} \times w \times \varepsilon_{\mathrm{adj}}
  \times (63/64)$, with $w = 7/80$
\item $\alpha^{-1} = 6\pi\,\kappa_{U(1)} = 137.03599985$ (sub-ppm)
\item No GUT-scale unification required; $\alpha$ predicted at low energy
\item Zero continuous free parameters
\end{itemize}

\end{tcolorbox}

%=============================================================================
\section{Identified Gaps and Future Work}
\label{sec:gaps}
%=============================================================================

\begin{gap}[No explicit $a_4$ evaluation on $\mathbb{C}P^2 \times S^3$]
The DFD corpus asserts that the gauge kinetic term is ``extracted from $a_4$''
but never evaluates the standard Seeley--DeWitt--Gilkey formula
on the specific geometry with Toeplitz truncation.
A complete derivation should evaluate the product-formula decomposition
$a_4(X) = \sum_{p+q=4} a_p(\mathbb{C}P^2) \cdot a_q(S^3)$ with the known
curvature invariants ($\mathbb{C}P^2$: K\"ahler--Einstein with $R = 24$;
$S^3$: constant curvature with $R = 6/r^2$).
\end{gap}

\begin{gap}[Connection to CCM traces not made explicit]
The DFD sources never explicitly compare the DFD trace coefficients
($w = 7/80$, $g_F = 8$, $\mathrm{Tr}(Y^2) = 10$) to the standard CCM traces
($c_1 = 10$, $c_2 = c_3 = 6$) or show how one framework maps to the other.
\end{gap}

\begin{gap}[$C_{\mathrm{loop}} = 1/8$ unverified]
The self-coupling relation $k_a = 3/(8\alpha)$ depends on $C_{\mathrm{loop}} = 1/8$,
described as ``plausible from heat kernel structure but requires explicit
verification.''  This should emerge from the $a_4$ computation on
$\mathbb{C}P^2 \times S^3$.
\end{gap}

\begin{gap}[Precise mapping between $w$ and $c_1$]
The hypercharge weighting $w = N_{\mathrm{species}}/(g_F \cdot \mathrm{Tr}(Y^2))
= 7/80$ plays the role of the CCM trace coefficient $c_1$, but the precise
relationship $c_1^{\mathrm{DFD}} = (\cdots) \times w$ has not been derived.
In particular, $g_F = 8$ should emerge from the KO-dimension of the finite
spectral triple ($2^3 = 8$ from $J \times \gamma \times \mathbb{C}$), and
$N_{\mathrm{species}} = 7$ should be derivable from the spectral triple structure.
\end{gap}

%=============================================================================
\section*{References}
%=============================================================================

\begin{enumerate}
\item[{[1]}] A.~H.~Chamseddine and A.~Connes, ``The Spectral Action Principle,''
  \textit{Commun.\ Math.\ Phys.}\ \textbf{186}, 731 (1997);
  arXiv:hep-th/9606001.

\item[{[2]}] A.~H.~Chamseddine, A.~Connes, and M.~Marcolli, ``Gravity and the
  Standard Model with Neutrino Mixing,'' \textit{Adv.\ Theor.\ Math.\ Phys.}\
  \textbf{11}, 991 (2007); arXiv:hep-th/0610241.

\item[{[3]}] A.~Devastato, M.~Kurkov, and F.~Lizzi, ``Spectral Noncommutative
  Geometry, Standard Model and all that,'' \textit{Int.\ J.\ Mod.\ Phys.\ A}\
  \textbf{34}, 1930010 (2019); arXiv:1906.09583.

\item[{[4]}] W.~D.~van Suijlekom, \textit{Noncommutative Geometry and Particle
  Physics}, 2nd ed.\ (Springer, 2024).

\item[{[5]}] G.~Alcock, ``Density Field Dynamics: A Complete Unified Theory,''
  v108, Zenodo (2026).  Section~8C, Appendices F, K, Z.

\item[{[6]}] G.~Alcock, ``Ab Initio Derivation of the Fine Structure Constant
  from Density Field Dynamics,'' (2025).
\end{enumerate}

\end{document}
